Pocklington's algorithm is a technique for solving congruences of the form $x^2 \equiv a \mod p$, where $x$ and $a$ are integers and $a$ is a quadratic residue modulo $p$. Given an integer $a$ and odd prime $p$ as input, the algorithm separates three cases for $p$, and then computes $x$ accordingly. We are interested in the case where  $p = 8m + 5$ for some $m\in\NN$, as we note that $x$ can be computed deterministically as follows:

\begin{figure}[H]
    \centering
    \begin{tabularx}{0.82\textwidth}{ p{0.1\textwidth} p{0.72\textwidth} }
   \rowcolor{Gray}
    \multicolumn{2}{c}{\textbf{Pocklington's algorithm for $p=8m+5$}}
    \\
    \textbf{Input:} & Prime $p \equiv 5 \mod 8$ and integer $a$ that is a quadratic residue modulo $p$ \\ 
    \rowcolor{Gray}
    \textbf{Output:} & Solution of the equation $x^2 \equiv a \pmod{p}$
    \\ 
    \textbf{Step 1:} &
            Write $p = 8m + 5$ with $m\in\NN$.
  \\
    \textbf{Step 2:} &
       If $a^{2m+1} \equiv 1 \pmod{p}$:
    \\ & $\qquad$ Return $x = \pm a^{m+1}$.
    \\ & If $a^{2m+1} \equiv -1 \pmod{p}$:
    \\ & $\qquad$ Write $y = \pm(4a)^{m+1}$ and return $x=\pm\frac{y}{2}$ if $y$ is even, and
    $x=\pm\frac{p+y}{2}$ if $y$ is odd.\\
   \hline
       \end{tabularx}
    \caption{Procedure to solve the congruence $x^2\equiv a \mod p$ for prime $p \equiv 5 \mod 8$.}
    \label{fig:pocklington}
\end{figure}

Let us briefly elaborate on the correctness of the above procedure. Given $p = 8m + 5$ with $m\in\NN$, following Fermat's little theorem, we have $x^{8m+4} \equiv 1 \mod p$. Since $x^2\equiv a \mod p$, we can rewrite it as $a^{4m+2} \equiv 1 \mod p$. Now, let us look at two separate cases for $a$:
\begin{enumerate}
    \item $a^{2m+1} \equiv 1 \mod p$ \\
    Then $a^{2m+2} \equiv a \mod p$, that is $(a^{m+1})^2 \equiv a \mod p$, hence $x = \pm a^{m+1}$.
    \item $a^{2m+1} \equiv -1 \mod p$ \\
    Note that 2 is a quadratic non-residue, so $4^{2m+1}\equiv -1 \mod p$, hence $4^{2m+1}a^{2m+1}\equiv 1 \mod p$. That is $(4a)^{2m+1}\equiv 1 \mod p$ and following the reasoning from the previous case $y = \pm(4a)^{m+1}$ is a solution of $y^2=4a$, hence $x=\pm\frac{y}{2}$, or if $y$ is odd, $x=\pm\frac{p+y}{2}$.
\end{enumerate}